Електронната търговия е сред най-динамично развиващите се области на съвременните информационни системи. Днешните онлайн магазини са комплексни приложения, които съчетават идентификация на потребители, представяне на продукти, управление на наличности, обработка на поръчки и устойчиво съхранение на данни. Предметната област на настоящата работа е продажба на компютърна и офис техника -- продукти, масово търсени от индивидуални и фирмени клиенти, и предполагащи структуриране по категории, проследяване на наличности и асоцииране с доставчици.

Основната цел на дипломната работа е да се проектира и разработи уеб базирана система за електронна търговия с компютърна и офис техника, която поддържа регистрация на потребители, управление на фирми и продукти, разглеждане на каталог, работа с количка, финализиране на поръчка и проследяване на история. За постигането ѝ са формулирани задачите: анализ на предметната област, избор на технологичен стек, проектиране на многослойна архитектура, моделиране на домейн обектите, реализация на регистрация и сесийно управление, изграждане на каталог, реализация на checkout и история, базова локализация и примерен REST интерфейс.

Проектът е реализиран като Java уеб приложение върху Spring Boot 2.6.2 и Java 11. Уеб слоят използва Spring MVC и Thymeleaf, persistence слоят -- Spring Data JPA и MySQL, а валидацията и сигурността -- Hibernate Validator, ModelMapper и Spring Security Crypto чрез \texttt{Pbkdf2PasswordEncoder}. Обхватът включва back-end логиката, HTML интерфейс, MySQL persistence и Maven build. Извън него остават реална платежна интеграция, Spring Security, Docker и разширено автоматизирано тестване.

Изложението е организирано в четири глави. Първата въвежда контекста на уеб-базираните платформи и теоретичните основи. Втората представя проектирането -- изисквания, архитектура и домейн модел. Третата описва програмната реализация на back-end, представящия слой и операционната среда. Четвъртата представя ръководство за използване и верификация. Заключението обобщава постигнатото, оценява решението и очертава насоки за бъдещо развитие.
